% from ega4-18-fr.tex
\oldpage[IV-4]{127}

On a par hypothèse $X=\Spec(\mathscr B)$, où $\mathscr B=f_*(\sh{O}_X)$ est une $\sh{O}_S$-Algèbre finie et localement libre. Posons $X'=X\times_S S'$, $f'=f_{(S')}$, $\mathscr B'=\mathscr B\otimes_{\sh{O}_S}\sh{O}_{S'}=f'_*(\sh{O}_{X'})$, de sorte que $X'=\Spec(\mathscr B')$; il y a alors une bijection canonique, fonctorielle en $S'$, de l'ensemble $\operatorname{Of}(X')$ sur l'ensemble $\operatorname{Id}(\mathscr B')$ des \emph{idempotents} de l'anneau $\Gamma(S',\mathscr B')=\Gamma(X',\sh{O}_{X'})$. En effet, en vertu de l'\emph{équivalence} de la catégorie des $S'$-schémas affines sur $S'$ et de la catégorie opposée de la catégorie des $\sh{O}_{S'}$-Algèbres quasi-cohérentes (II, 1.2.7 et 1.3.1), il y a correspondance biunivoque canonique entre les décompositions de $X'$ en \emph{somme} $X'_1\amalg X'_2$ de deux sous-préschémas induits sur des ouverts de $X'$ et les décompositions de $\mathscr B'$ en \emph{composée directe} de deux Idéaux $\mathscr B'_1$, $\mathscr B'_2$; ces dernières à leur tour forment un ensemble en correspondance biunivoque canonique (et fonctorielle en $S'$) avec $\operatorname{Id}(\mathscr B')$. Il suffit donc de prouver le lemme pour le foncteur $S'\longmapsto\operatorname{Id}(\mathscr B')$.

Pour cela, nous allons montrer qu'il existe un Idéal de type fini $\mathscr J$ de $\mathbf S_{\sh{O}_S}(\check{\mathscr B})$ tel que $\operatorname{Id}(\mathscr B')$ soit de la forme $\operatorname{Al}(S',\mathscr B',\mathscr J\otimes_{\sh{O}_S}\sh{O}_{S'})$ (18.5.2). Notons à cet effet que puisque $\mathscr B$ est une $\sh{O}_S$-Algèbre, son image réciproque $\mathscr B_{(\mathbf V(\check{\mathscr B}))}$ est une $\sh{O}_{(\mathbf V(\check{\mathscr B}))}$-Algèbre, et l'on peut donc former le \emph{carré} $c^2$, dans cette Algèbre, de la section canonique $c$ (18.5.1); elle correspond canoniquement à un homomorphisme $u^{(2)}:\mathbf S_{\sh{O}_S}(\check{\mathscr B})\to\mathbf S_{\sh{O}_S}(\check{\mathscr B})\otimes_{\sh{O}_S}\mathscr B$ de $\mathbf S_{\sh{O}_S}(\check{\mathscr B})$-Modules, et on vérifie aussitôt que pour un ouvert affine $W$ de $S$, avec les notations de (18.5.1), $u^{(2)}$ correspond à l'homomorphisme de $C$-modules tel que
\[
  u^{(2)}(1)=\sum_k\Bigl(\sum_{i,j}c_{ijk}T_iT_j\Bigr)\otimes e_k,
\]
où $(c_{ijk})$ est la table de multiplication de l'algèbre $\Gamma(W,\mathscr B)$. Montrons que l'Idéal $\mathscr J$ de $\mathbf S_{\sh{O}_S}(\check{\mathscr B})$ engendré par le \emph{noyau} de l'homomorphisme $u-u^{(2)}$ répond à la question. Il suffit en effet de noter que l'idéal $\Gamma(W,\mathscr J)$ de $C$ est engendré par les polynômes $P_k(T_1,\ldots,T_n)=T_k-\sum_{i,j}c_{ijk}T_iT_j$ et que les idempotents de $\Gamma(W,\mathscr B)$ sont précisément les éléments $\sum_i t_i e_i$ de cette algèbre tels que $(t_1,\ldots,t_n)$ annule tous les polynômes $P_k$ ($1\leq i\leq n$). On déduit donc de (18.5.2) que le $S$-préschéma affine de présentation finie $\mathbf{Of}(X)=\mathbf{Al}(\mathscr B,\mathscr J)$ représente bien le foncteur (18.5.3.1). Il reste à prouver que $\mathbf{Of}(X)$ est étale sur $S$, ou, ce qui revient au même, qu'il est \emph{formellement étale} sur $S$. Mais si $S'$ est un $S$-préschéma, $S'_0$ un sous-préschéma fermé de $S'$ défini par un Idéal localement nilpotent de $\sh{O}_{S'}$ (et ayant donc même espace sous-jacent que $S'$), il est clair que $X'=X\times_S S'$ et $X'_0=X\times_S S'_0$ ont même espace sous-jacent, donc l'application canonique $\operatorname{Of}(X')\to\operatorname{Of}(X'_0)$ est \emph{bijective}, ce qui achève la démonstration (17.1.1).

\begin{proposition}[18.5.4]
\label{IV.18.5.4}
Soient $S$ un préschéma, $S_0$ un sous-préschéma fermé de $S$; considérons les propriétés suivantes:
\begin{enumerate}
  \item[a)] Pour tout morphisme fini $g:S'\to S$, l'application canonique
  \begin{equation}
    \operatorname{Of}(S')\longrightarrow\operatorname{Of}(S'\times_S S_0)
    \tag{18.5.4.1}\label{IV.18.5.4.1}
  \end{equation}
  (cf. (18.5.3)) est bijective.
  \item[a')] Pour tout morphisme fini et localement libre $g:S'\to S$, l'application (18.5.4.1) est bijective.

